% Stage-1 commitment optimizes by simulating Stage-2 recourse over sampled
% populations: users enter first as model draws, then as real arrivals.
% Single-column figure for the approach section.
\begin{figure}[t]
\centering
\begin{tikzpicture}[
  font=\footnotesize,
  every node/.style={transform shape},
  stg/.style={rectangle, rounded corners=3pt, draw=black!55, thick,
              align=left, inner sep=6pt, text width=0.40\textwidth},
  dep/.style={-{Latex[length=2mm]}, very thick, black!75},
  lo/.style={-{Latex[length=1.8mm]}, thick, black!55, densely dashed}
]
\node[stg, fill=blue!5] (s1) {%
{\bfseries Stage 1: commit $\FSL_d$ (end of $d{-}1$)}\\[2pt]
score candidate commitments by \emph{simulating the day below} on
populations sampled from $\psihat$; keep the $\varepsilon$-feasible set,
pick depth $\lambda$.\\[2pt]
{\scriptsize\itshape users enter here only as sampled draws}};

\node[stg, fill=teal!6, below=0.9cm of s1] (s2) {%
{\bfseries Stage 2: the operating day (recourse)}\\[2pt]
real arrivals negotiate under the committed $\FSL_d$; charging banks,
delivers, and settles on realized quantities.\\[2pt]
{\scriptsize\itshape the same machinery Stage 1 simulated, now on real users}};

\draw[dep] (s1) -- (s2) node[midway, right=2pt, font=\scriptsize\itshape, text=black!60]{commit, then face the day};

\draw[lo] (s1.west) ++(-0.05,0) arc[start angle=-60, end angle=240, radius=0.28]
  node[left=6pt, font=\scriptsize\itshape, text=black!60, align=right]{simulate\\per sample};
\end{tikzpicture}
\caption{The commitment is chosen by rehearsing its own consequences: Stage 1
scores each candidate $\FSL_d$ by running the Stage-2 negotiation and
charging machinery on populations sampled from the forecast, then commits;
on the day itself the identical machinery runs on the real arrivals. Users
appear twice, first as model draws inside the rehearsal, then as arriving
negotiators, which is what makes the commitment feasible before its
resource exists.}
\label{fig:twostage}
\end{figure}
